\documentclass[10pt,a4paper]{amsart}
\usepackage[margin=19mm]{geometry}
\usepackage{amsmath,amssymb,mathtools}
\usepackage[scr=rsfs,cal=euler]{mathalfa}
\usepackage{xfrac}
\usepackage[unicode,hidelinks,bookmarks=false]{hyperref}
\newcommand{\ie}{i.e.\ }
\newcommand{\cf}{cf.\ }
\newcommand{\resp}{respectively\ }
\newcommand{\Subsection}{\subsection}
\providecommand{\nobreakdash}{\nobreak\hbox{-}}
\setlength{\emergencystretch}{2em}
\multlinegap0pt
\usepackage{amssymb}
\usepackage{enumerate}
\usepackage{caption}
\usepackage[all]{xy}
\usepackage{tikz}
\newtheorem{lemma}{Lemma}
\newcommand{\Q}{\mathbb{Q}}
\title{Non-commutative crepant resolutions for (almost) simplicial toric algebras}
\author{Aimeric Malter, Artan Sheshmani}
\date{Source: arXiv:2602.21802v1 (CC BY 4.0)}
\begin{document}
\maketitle

\noindent\emph{Source and licence.} Non-commutative crepant resolutions for (almost) simplicial toric algebras. Version arXiv:2602.21802v1 (CC BY 4.0), distributed by arXiv under the Creative Commons Attribution 4.0 International licence, http://creativecommons.org/licenses/by/4.0/, as recorded in the arXiv licence metadata for this exact version and shown on the abstract page https://arxiv.org/abs/2602.21802v1. Full licence notice: \texttt{licenses/arxiv-2602.21802-CC-BY-4.0.txt}.\par\smallskip

\noindent\textit{Editorial scope.} Counted unit: the article's lemma Lem:collectionsrandalpha (source0.tex lines 438--441). The article's complete original proof of it is delivered with it (source0.tex lines 442--447, one \texttt{proof} environment of 6 lines). No mathematics is written, repaired or re-derived: the statement and the proof are the article's own lines, in the article's own notation, and no retained line is substituted or re-flowed.  No further source-local context is needed: every cross-reference of every retained line resolves inside the two delivered spans.  Nothing was omitted from any retained span, so the package carries no omission record.  No retained line contains a citation, so the wrapper carries no bibliography and no bibliography entry.  No retained line contains a figure environment, an image-inclusion command or drawing code, so no image is delivered and the package declares no asset dependency.  Editorial formatting only, in addition: an AMS \texttt{amsart} preamble that restates the article's own declarations and macros used by the retained lines, namely the environments \texttt{lemma} on separate counters, together with the standard packages those lines need.  The retained environments print in this document as Lemma 1 in order of appearance, whereas the article prints the same items with the numbering of its own full document.  The complete native source of the article as distributed in the arXiv e-print for this exact version is bundled unchanged as \texttt{sources/195282/source0.tex}.
\begin{lemma}
    \label{Lem:collectionsrandalpha}
    There exists a collection of positive numbers $r_i$ such that $\sum_{i=1}^{n+3}r_i v_i=0$ and $\sum_{i=1}^{n+3}r_i=1$. There also exists a, unique up to scaling, collection of rational numbers $\alpha_i$ such that $\sum_{i=1}^{n+3}\alpha_i v_i=0$ and $\sum_{i=1}^{n+3}\alpha_i=0$. Furthermore, $\alpha_{n+3}=0$. 
\end{lemma}

\begin{proof}
Begin by noting that $0$ is in the interior of $Q$, implying the existence of the collection $r_i$ directly. 
To find the collection $\alpha_i$, note that the fan $\Sigma$ is complete and its primitive vertices $v_i$ generate $N\otimes \Q$, and so the space of linear relations has dimension two. The collection of $r_i$ above has $\sum_{i=1}^{n+3}r_i>0$, so the condition $\sum_{i=1}^{n+3}\alpha_i=0$ cuts out a dimension one subspace of the linear relations. We claim now that we can pick this collection such that $\alpha_1,\alpha_2>0$, $\alpha_3,\dots,\alpha_{n+2}<0$ and $\alpha_{n+3}=0$. Note that the line connecting $w_1$ and $w_2$ intersects the hyperplane $H$ in a unique point $z\in \operatorname{Relint}(H)$.

Thus, there are positive coefficients $\alpha_1, \alpha_2>0$ with $\alpha_1w_1+\alpha_2 w_2=z$ and $\alpha_1+\alpha_2=1$ as well as coefficients $\beta_i>0$ such that $\sum_{i=3}^{n+2}\beta_i w_i=z$ and $\sum_{i=3}^{n+2}\beta_i=1$. The image $\tilde{z}$ of the point $z$ in the face of $Q$ obtained by placing $P$ has the same property, i.e. $\tilde{z}=\alpha_1 v_1+\alpha_2 v_2=\sum_{i=3}^{n+2}\beta_i v_i$. Write $\alpha_i=-\beta_i$ for $3\le i\le n+2$ to obtain $\sum_{i=1}^{n+2}\alpha_i v_i=0$ and $\sum_{i=1}^{n+2}\alpha_i=0$. Adding $\alpha_{n+3}=0$ we obtain a collection of $\alpha_i$ with $\sum_{i=1}^{n+3}\alpha_i v_i=0$ and $\sum_{i=1}^{n+3}\alpha_i=0$, and so by uniqueness up to scaling our claim holds.
\end{proof}

\end{document}
